\scanpage{37}
\begin{center}- 44 -\end{center}
\hfill n$^\circ$ 272

$\Gamma_{\unsure{\omega|w}}/\Gamma_v$ sont finis, le groupe $\Gamma_{v'}/\Gamma_v$ est un groupe de torsion. On raisonne de même pour montrer que $L_{A'}/K_A$ est algébrique.

\uline{Corollaire} 1 - \uline{Le rang de} $v'$ \uline{est égal à celui de} $v$.

Posons pour simplifier les notations, $G=\Gamma_v$, $G'=\Gamma_{v'}$, et soit $S$ (resp. $S'$) l'ensemble des sous-groupes isolés de $G$ (resp. $G'$). Nous allons montrer que l'application $H'\longrightarrow H'\cap G$ est une bijection de $S'$ sur $S$, ce qui établira évidemment que $v$ et $v'$ ont même rang. Tout d'abord, il est clair que $H'\in S'\Longrightarrow H'\cap G\in S$ ; si de plus $H'_1\subset H'_2$ sont deux sous-groupes isolés de $G'$ tels que $H'_1\cap G=H'_2\cap G$, le groupe quotient $H'_2/H'_1$ est isomorphe à un sous-groupe de $G'/G$ ; comme le premier groupe est un groupe totalement ordonné et que le \unsure{sedond|second} est un groupe de torsion d'après la prop.2, il s'ensuit que $H'_2/H'_1=0$, et l'application considérée est injective. Enfin, si $H$ est un sous-groupe isolé de $G$, soit $H'\subset G'$ l'ensemble des $x\in G'$ tels qu'il existe $h\in H$ avec $-h\leqslant x\leqslant h$ ; on vérifie immédiatement que $H'$ est un sous-groupe isolé de $G'$ tel que $H'\cap G=H$, ce qui achève la démonstration.

[Est-il nécessaire de dégager \unsure{lo|le} ou les lemmes sur les groupes ordonnés que l'on vient \unsure{en'fait|en fait} de démontrer ?]

\uline{Corollaire} 2 - \uline{Pour que} $v'$ \uline{soit triviale} (resp. \uline{de rang} 1), \uline{il faut et il suffit que} $v$ \uline{soit triviale} (resp. \uline{de rang} 1).

Ce sont les cas particuliers du corollaire précédent correspondant aux valeurs 0 et 1 du rang.

\uline{Corollaire} 3 - \uline{Supposons que} $L$ \uline{soit une extension finie de} $K$. \uline{Alors, pour que} $v'$ \uline{soit discrète, il faut et il suffit que} $v$ \uline{le soit}.

Si $v'$ est discrète, le groupe $\Gamma_v$ est isomorphe à un sous-groupe de $\underline{\underline Z}$, non réduit à 0 d'après le cor.2 donc $v$ est discrète. Réciproquement, si $v$ est discrète, la prop.1 montre que $\Gamma_{v'}/\Gamma_v$ est un groupe fini, donc $\Gamma_{v'}$
\mspage{36}
